\section{Background: from surrogate selection to field ensembles}
\label{sec:related}
Engineering studies already show why choosing one surrogate can be difficult
and why combining several can help. \citet{goel2007} weight surrogates using
error estimates, \citet{acar2009} and \citet{stromberg2021} optimize metamodel weights, and
\citet{viana2009} compare cross validation selection, weighted prediction, and
library subsets. For functional outputs, \citet{brunel2025} compare more than
a dozen multifidelity surrogates in a common framework. Their findings connect
model performance to fidelity correlation, training size, and residual
complexity. These studies establish both the ensemble principle and the need
for comparisons across problems. We extend this empirical line of work to a
surrogate library for fields and measure combination gains under small,
explicitly counted fitting budgets.

The weighting rules follow stacked generalization \citep{wolpert1992}, stacked
regression \citep{breiman1996}, and Super Learner \citep{superlearner2007}.
Greedy ensemble selection \citep{caruana2004}, auto sklearn
\citep{feurer2015,feurer2022}, and AutoGluon \citep{erickson2020} automate
selection from broader libraries. Related multifidelity workflows automate
other parts of the process: EL MFS combines low fidelity models inside
hierarchical fusion \citep{wang2024elmfs}, AQBMF screens low fidelity information
and surrogate candidates \citep{lee2026aqbmf}, and the MAESTRO preprint screens
families by leave one out validation \citep{kim2026maestro}. Our ensemble fitting
stage acts on completed target field predictions from different fidelity
strategies. We evaluate that stage and its library, without a direct comparison
against these complete engineering automation systems.

The library draws on Fourier neural operators (FNO) \citep{li2021fno}, DeepONet
\citep{lu2021}, wavelet neural operators (WNO) \citep{tripura2022wno}, and Transolver
\citep{wu2024transolver}. Multifidelity operator learning also includes
DeepONet \citep{howard2023}, MF WNO \citep{thakur2022mf}, and FNO for carbon
storage \citep{tang2023}. MFRNP uses decoded coarse outputs in residual prediction
\citep{niu2024}, while FIRE uses coarse predictive distributions
\citep{yu2026fire}. Appendix~\ref{app:adaptations} identifies the components we
adapt and the departures from their original formulations. HyperNOs addresses
automated operator training \citep{ghiotto2025}, and LANE SI combines sea ice
fields \citep{borisova2024}. PDEBench and The Well supply broader physical
benchmarks \citep{takamoto2022,ohana2024}. Our study starts from supplied training
data and fits ensemble weights, leaving simulation acquisition to approaches
such as active multi resolution learning \citep{li2024active}.
